% ===========================================================================
\section{最大公因数与扩展欧几里得算法}\label{sec:gcd}

本节证明上节预告的等式 $\gcd(a,b)=\gcdp(a,b)$。证明分三步，分别给出三个量
之间的不等式，最后夹出相等。记 $\mathrm{GCD}(a,b)$ 为下面算法的输出。

\block{第一步：$\gcd(a,b)\le\gcdp(a,b)$}
令 $d=\gcd(a,b)$，并设 $\gcdp(a,b)=am'+bn'$。因为 $d\mid a$ 且 $d\mid b$，
显见 $d\mid am'+bn'$，即 $\gcd(a,b)\mid\gcdp(a,b)$，于是
\begin{equation}
  \gcd(a,b)\le\gcdp(a,b). \tag{1}
\end{equation}

\block{第二步：扩展欧几里得算法}
\textbf{算法 $\mathrm{GCD}$（扩展欧几里得算法）。}
输入 $0\le a\le b$；输出 $(d,m,n)$，满足 $d=am+bn$。
\begin{itemize}[leftmargin=1.6em,itemsep=1pt]
  \item 若 $a=0$，输出 $(b,\,0,\,1)$（此时 $d=b=0\cdot a+1\cdot b$）；
  \item 否则做带余除法 $b=aq+r$（$0\le r<a$），递归计算
        $(d,m,n)=\mathrm{GCD}(r,a)$，输出 $(d,\;n-qm,\;m)$。
\end{itemize}
递归的正确性由
\[
  d=rm+an=(b-aq)m+an=a(n-qm)+bm
\]
保证；第一参数从 $a$ 严格降为 $r$，故算法终止。由于输出 $d$ 是 $a,b$ 的
正整系数线性组合，由 $\gcdp$ 的最小性得
\begin{equation}
  \gcdp(a,b)\le\mathrm{GCD}(a,b). \tag{2}
\end{equation}

\block{第三步：$\mathrm{GCD}(a,b)\le\gcd(a,b)$}
对递归深度归纳可证 $\mathrm{GCD}(a,b)\mid a$ 且 $\mathrm{GCD}(a,b)\mid b$：
基例 $a=0$ 时输出 $b$，显然整除 $0$ 与 $b$；递归步中 $d\mid r$、$d\mid a$，
而 $b=aq+r$，故 $d\mid b$。于是 $\mathrm{GCD}(a,b)$ 是 $a,b$ 的公因子，
不超过最大公因子：
\begin{equation}
  \mathrm{GCD}(a,b)\le\gcd(a,b). \tag{3}
\end{equation}

\begin{theorem}[Bézout 恒等式]\label{thm:bezout}
设 $a,b$ 不全为 $0$，则
\[
  \gcd(a,b)=\gcdp(a,b)=\min\{am+bn>0:m,n\in\Z\}.
\]
特别地，存在整数 $m,n$ 使 $\gcd(a,b)=am+bn$，且这样的系数可由扩展欧几里得
算法实际算出。
\end{theorem}

\begin{proof}
结合 (1)(2)(3) 得
$\gcdp(a,b)\le\mathrm{GCD}(a,b)\le\gcd(a,b)\le\gcdp(a,b)$，故三者相等。
\end{proof}

\begin{corollary}[互素判据]\label{cor:coprime-crit}
$\gcd(a,b)=1$ 当且仅当存在整数 $m,n$ 使 $am+bn=1$。
\end{corollary}

\begin{definition}[互素]
若 $\gcd(a,b)=1$，称 $a$ 与 $b$ \keyword{互素}（coprime），记作 $a\cop b$。
\end{definition}

\begin{lemma}[互素的乘积封闭性]\label{lem:cop-product}
若 $\gcd(a,b)=1$ 且 $\gcd(a,c)=1$，则 $\gcd(a,bc)=1$。
\end{lemma}

\begin{proof}
由推论 \ref{cor:coprime-crit}，存在 $m,n,m',n'\in\Z$ 使
\[
  am+bn=1,\qquad am'+cn'=1.
\]
两式相乘（写成 $bn=1-am$、$cn'=1-am'$ 的形式）得
\[
  bc\cdot nn'=(1-am)(1-am')=1-a(m+m'-amm'),
\]
即
\[
  a\,(m+m'-amm')+bc\,(nn')=1.
\]
再用一次推论 \ref{cor:coprime-crit}，得 $\gcd(a,bc)=1$。
\end{proof}

\begin{remark}
这条引理是中国剩余定理唯一性论证的关键：由它归纳可得，若
$m_1,\ldots,m_t$ 两两互素且每个 $m_i$ 都整除 $x$，则乘积
$m_1\cdots m_t$ 也整除 $x$。
\end{remark}
